% Subsection 4.4.2: Few-shot CoT 注入

Few-shot CoT（CoT-FS($k$)）在 CoT-ZS 基础上的要点如下：
\begin{itemize}
    \item \textbf{示例形态}：每条为 \texttt{(题目, 推理, 代码)}，在 chat 中编码为 \textbf{User(CoT 版题指令)} $\rightarrow$ \textbf{Assistant(推理 + fenced Python)}；
    \item \textbf{两阶段对称回放}：Stage-1 在算法教师 system 之后插入 $k$ 组示范，再接目标题 User；Stage-2 在程序员 system 之后用\textbf{同一组}示例以\textbf{相同顺序}再回放一遍，最后接 \texttt{Reasoning:} + 截断 Stage-1 + 目标题的综合 User；
    \item \textbf{实现对应}：\texttt{eval\_lora\_cot\_passk.py}；训练侧 few-shot 交替结构见 \texttt{train\_cot.py}；
    \item \textbf{实验编号}：H3、H6、H10；
    \item \textbf{选池原则}：按强标签检索，训练与评测共享同一评分函数，以保证示例与目标题尽量同构。
\end{itemize}

\subsubsection{候选评分与选池}
对每条目标样本 $q$，从 3.3.2 节定义的强标签集合派生 \texttt{strong\_tags}，在候选集合（训练集或评测集扣除 $q$）上按下式打分：
\[
\mathrm{score}(q, q') = 2 \cdot \bigl|\,\mathrm{strong\_tags}(q) \cap \mathrm{strong\_tags}(q')\,\bigr|
                       + \bigl|\,\mathrm{tags}(q) \cap \mathrm{tags}(q')\,\bigr|.
\]
\begin{itemize}
    \item \textbf{加权含义}：强标签命中权重 2，普通标签权重 1，兼顾「算法骨架一致」与「题型相近」；
    \item \textbf{候选池规模}：按分数降序取前 $\max(k, 4k)$ 条；
    \item \textbf{池内抽样}：再随机抽样，避免总是固定命中极少数高分题。
\end{itemize}

\subsubsection{训练/评测两侧的随机性差异}
\begin{itemize}
    \item \textbf{共同点}：两侧共用同一评分函数与同一强标签集合，差异仅在随机采样；
    \item \textbf{训练侧}：\texttt{random.sample(pool, k)}，依赖 \texttt{seed=42} 的全局 RNG，每个 epoch 的示例在统计上来自同一分布；
    \item \textbf{评测侧}：\texttt{seed = 2026 + problem\_idx * 1000 + sample\_idx}，同一 \texttt{(题目, 采样编号)} 的示例选择完全可复现，便于跨实验、跨运行做差分归因；
    \item \textbf{合取含义}：训练与评测在「评分 + 候选池构造」上统计一致；评测侧因固定 seed，任意第三方可按题号与采样编号逐条复核。
\end{itemize}

\subsubsection{示例拼接与预算分配}
\begin{itemize}
    \item \textbf{多轮形态}：$k$ 条示例在 Stage-1 与 Stage-2 中均为「\textbf{User（CoT 版题目指令）} $\rightarrow$ \textbf{Assistant（推理 + fenced Python）}」，置于目标轮次之前（与 4.1 节（1b）一致），再接目标题 User 或 \texttt{Reasoning:} 综合 User；
    \item \textbf{三类长度预算}（两阶段 prompt 共享）：
    \begin{itemize}
        \item 示例推理保持原文，不在此阶段额外截断（长度由数据生成阶段约束）；
        \item 目标题的 Stage-1 输出仍受 \texttt{MAX\_REASONING\_CHARS} 与 \texttt{MAX\_REASONING\_PROMPT\_TOKENS} 二级截断；
        \item \texttt{VLLM\_MAX\_MODEL\_LEN=8192} 覆盖「全部示例 + 目标题 + 两阶段生成」。
    \end{itemize}
\end{itemize}

\subsubsection{训练/推理 $k$ 匹配的强约束}
\begin{itemize}
    \item \textbf{规则来源}：4.1 节合法性约束；训练侧为 \texttt{LoRA\_cot\_fs($k$)} 时，推理 \texttt{--few-shot-k} 必须同为 $k$；
    \item \textbf{脚本行为}：不匹配时评测脚本显式退出（H10 对应 \texttt{FS\_K}）；
    \item \textbf{归因意义}：保证「训练见 $k$ 例、评测亦见 $k$ 例」，是 few-shot 净效应可比的前提；否则 train/test 在上下文长度与示例密度上错位，会混淆训练侧与推理侧贡献。
\end{itemize}